% GENERATED FILE. Edit canonical lessons, then run npm run generate:books.
% canonical-source-hash: fnv1a64:c79091a9e1aa70b3
% canonical-lessons: FR-C140-deja, FR-C140-recemment, FR-C140-avant-hier, FR-C140-autrefois, FR-C140-se-produire

\chapter{Already, Recently, the Day before Yesterday, Formerly, To Happen}
\label{ch:fr-already-recently-the-day-before-yesterday-formerly-to-happen}
\hlchaptermodality{140}

\begin{chapteropening}
\textbf{By the end of this chapter:} \emph{I can say what happened, and when: already, recently, the day before yesterday.}

Complete the last lesson of chapter 140: I can say what happened, and when: already, recently, the day before yesterday.
\end{chapteropening}

\section[déjà]{déjà — already}

\label{lesson:FR-C140-deja}



\begin{quote}
Before the new one: say the French for which one, then the French for anything at all.
\end{quote}

\subsection*{You'll want to know: déjà}
\textbf{déjà} — \textquotedblleft{}already\textquotedblright{}. It sits between \textbf{avoir} and the participle: \textbf{J'ai déjà parlé.} — \textquotedblleft{}I have already spoken.\textquotedblright{}

\begin{grammarlens}[title={What you've built}]
One more word for this chapter.
\end{grammarlens}

\subsection*{Your turn}
\begin{itemize}
\raggedright
  \item \emph{Say it:} \emph{déjà}
  \item \emph{Say it:} \emph{déjà}, once more
  \item \emph{Say it:} say \emph{déjà}
  \item \emph{Recall:} say the French for which one, then the French for anything at all, then say \emph{déjà} again
\end{itemize}

\begin{tcolorbox}[breakable,colback=teal!4,colframe=teal!35!black,title={Before you move on}]
What does déjà mean? (\textquotedblleft{}Already\textquotedblright{}.) Say it once more.
\end{tcolorbox}
\section[récemment]{récemment — recently}

\label{lesson:FR-C140-recemment}



\begin{quote}
Before the new one: say the French for anything at all, then the French for already.
\end{quote}

\subsection*{You'll want to know: récemment}
\textbf{récemment} — \textquotedblleft{}recently\textquotedblright{}. \textbf{Récemment, j'ai beaucoup lu.} — \textquotedblleft{}Recently I have read a lot.\textquotedblright{} The ending \textbf{-ment} makes an adverb, the way English uses -ly.

\begin{grammarlens}[title={What you've built}]
One more word for this chapter.
\end{grammarlens}

\subsection*{Your turn}
\begin{itemize}
\raggedright
  \item \emph{Say it:} \emph{récemment}
  \item \emph{Say it:} \emph{récemment}, once more
  \item \emph{Say it:} say \emph{récemment}
  \item \emph{Recall:} say the French for anything at all, then the French for already, then say \emph{récemment} again
\end{itemize}

\begin{tcolorbox}[breakable,colback=teal!4,colframe=teal!35!black,title={Before you move on}]
What does récemment mean? (\textquotedblleft{}Recently\textquotedblright{}.) Say it once more.
\end{tcolorbox}
\section[avant-hier]{avant-hier — the day before yesterday}

\label{lesson:FR-C140-avant-hier}



\begin{quote}
Before the new one: say the French for already, then the French for recently.
\end{quote}

\subsection*{You'll want to know: avant-hier}
\textbf{avant-hier} — \textquotedblleft{}the day before yesterday\textquotedblright{}. One day further back than \textbf{hier}: \textbf{Avant-hier, j'ai écrit une lettre.} — \textquotedblleft{}The day before yesterday I wrote a letter.\textquotedblright{}

\begin{grammarlens}[title={What you've built}]
One more word for this chapter.
\end{grammarlens}

\subsection*{Your turn}
\begin{itemize}
\raggedright
  \item \emph{Say it:} \emph{avant-hier}
  \item \emph{Say it:} \emph{avant-hier}, once more
  \item \emph{Say it:} say \emph{avant-hier}
  \item \emph{Recall:} say the French for already, then the French for recently, then say \emph{avant-hier} again
\end{itemize}

\begin{tcolorbox}[breakable,colback=teal!4,colframe=teal!35!black,title={Before you move on}]
What does avant-hier mean? (\textquotedblleft{}The day before yesterday\textquotedblright{}.) Say it once more.
\end{tcolorbox}
\section[autrefois]{autrefois — formerly, in the old days}

\label{lesson:FR-C140-autrefois}



\begin{quote}
Before the new one: say the French for recently, then the French for the day before yesterday.
\end{quote}

\subsection*{You'll want to know: autrefois}
\textbf{autrefois} — \textquotedblleft{}formerly, in the old days\textquotedblright{}. It points to a time long past rather than a single day: \textbf{Autrefois, le magasin était ici.} — \textquotedblleft{}In the old days the shop was here.\textquotedblright{}

\begin{grammarlens}[title={What you've built}]
One more word for this chapter.
\end{grammarlens}

\subsection*{Your turn}
\begin{itemize}
\raggedright
  \item \emph{Say it:} \emph{autrefois}
  \item \emph{Say it:} \emph{autrefois}, once more
  \item \emph{Say it:} say \emph{autrefois}
  \item \emph{Recall:} say the French for recently, then the French for the day before yesterday, then say \emph{autrefois} again
\end{itemize}

\begin{tcolorbox}[breakable,colback=teal!4,colframe=teal!35!black,title={Before you move on}]
What does autrefois mean? (\textquotedblleft{}Formerly, in the old days\textquotedblright{}.) Say it once more.
\end{tcolorbox}
\section[se produire]{se produire — to happen, to occur}

\label{lesson:FR-C140-se-produire}



\begin{quote}
Before the new one: say the French for the day before yesterday, then the French for formerly.
\end{quote}

\subsection*{You'll want to know: se produire}
\textbf{se produire} — \textquotedblleft{}to happen, to occur\textquotedblright{}. Its past takes \textbf{être}, like \textbf{aller}: \textbf{Qu'est-ce qui s'est produit ?} — \textquotedblleft{}What happened?\textquotedblright{}

\begin{grammarlens}[title={What you've built}]
One more word for this chapter.
\end{grammarlens}

\subsection*{Your turn}
\begin{itemize}
\raggedright
  \item \emph{Say it:} \emph{se produire}
  \item \emph{Say it:} \emph{se produire}, once more
  \item \emph{Say it:} say \emph{se produire}
  \item \emph{Recall:} say the French for the day before yesterday, then the French for formerly, then say \emph{se produire} again
\end{itemize}

\begin{tcolorbox}[breakable,colback=teal!4,colframe=teal!35!black,title={Before you move on}]
What does se produire mean? (\textquotedblleft{}To happen, to occur\textquotedblright{}.) Say it once more.
\end{tcolorbox}
